% ============================================================
\chapter{Algorithme de Gibbs}
\label{ch:gibbs}
% ============================================================

En 1984, Stuart et Donald Geman publient un article sur la restauration d'images en introduisant un algorithme MCMC particuli\`erement \'el\'egant~: l'\emph{\'echantillonneur de Gibbs}. Au lieu de proposer des mouvements dans l'espace complet des param\`etres (comme Metropolis--Hastings), l'algorithme met \`a jour chaque composante s\'epar\'ement, en tirant directement de sa \emph{loi conditionnelle compl\`ete} --- la distribution de cette composante conditionnellement \`a toutes les autres. Le g\'enie du proc\'ed\'e est que le taux d'acceptation est toujours de $100\,\%$, \'eliminant le besoin de calibrer une distribution de proposition. Alan Gelfand et Adrian Smith (1990) montrent ensuite que l'\'echantillonneur de Gibbs est applicable \`a une vaste classe de mod\`eles bay\'esiens hi\'erarchiques, d\'eclenchant la r\'evolution MCMC en statistique.

\begin{intuition}[title=\textbf{Id\'ee centrale}]
L'\'echantillonneur de Gibbs est un cas particulier de
Metropolis--Hastings o\`u chaque composante du param\`etre est mise \`a
jour en tirant directement depuis sa \textbf{loi conditionnelle compl\`ete}.
Le taux d'acceptation est toujours $1$, ce qui \'elimine le besoin de
calibration de la proposition.
\end{intuition}

% ------------------------------------------------------------
\section{Lois conditionnelles compl\`etes}
% ------------------------------------------------------------

\begin{definition}[Loi conditionnelle compl\`ete]
Soit $\theta = (\theta_1, \ldots, \theta_d)$. La \textbf{loi
conditionnelle compl\`ete} de $\theta_k$ est~:
\[
  \pi(\theta_k \mid \theta_{-k}, x)
  \propto \pi(\theta_1, \ldots, \theta_d \mid x)
  \quad \text{en tant que fonction de } \theta_k,
\]
o\`u $\theta_{-k} = (\theta_1, \ldots, \theta_{k-1}, \theta_{k+1}, \ldots, \theta_d)$.
\end{definition}

\begin{remarque}
Pour identifier la conditionnelle compl\`ete, on fixe toutes les
composantes sauf $\theta_k$ dans le posterior joint et on reconna\^it
la loi en $\theta_k$.
\end{remarque}

% ------------------------------------------------------------
\section{Algorithme de Gibbs}
% ------------------------------------------------------------

\begin{algorithme}[title=\textbf{\'Echantillonneur de Gibbs}]
\begin{enumerate}
  \item Initialiser $\theta^{(0)} = (\theta_1^{(0)}, \ldots, \theta_d^{(0)})$.
  \item Pour $t = 0, 1, \ldots, T-1$~:
  \begin{enumerate}
    \item Tirer $\theta_1^{(t+1)} \sim \pi(\theta_1 \mid \theta_2^{(t)}, \ldots, \theta_d^{(t)}, x)$.
    \item Tirer $\theta_2^{(t+1)} \sim \pi(\theta_2 \mid \theta_1^{(t+1)}, \theta_3^{(t)}, \ldots, \theta_d^{(t)}, x)$.
    \item[$\vdots$]
    \item Tirer $\theta_d^{(t+1)} \sim \pi(\theta_d \mid \theta_1^{(t+1)}, \ldots, \theta_{d-1}^{(t+1)}, x)$.
  \end{enumerate}
\end{enumerate}
\end{algorithme}

\begin{theoreme}[Gibbs comme cas particulier de MH]
L'\'echantillonneur de Gibbs est un algorithme de Metropolis--Hastings
o\`u la proposition pour $\theta_k$ est
$q_k(\theta_k^* \mid \theta_{-k}) = \pi(\theta_k^* \mid \theta_{-k}, x)$.
Le ratio d'acceptation vaut toujours~$1$.
\end{theoreme}

\begin{proof}
Le ratio d'acceptation pour la mise \`a jour de $\theta_k$ est~:
\begin{align*}
  \alpha &= \frac{\pi(\theta_k^*, \theta_{-k} \mid x)\,\pi(\theta_k^{(t)} \mid \theta_{-k}, x)}
                 {\pi(\theta_k^{(t)}, \theta_{-k} \mid x)\,\pi(\theta_k^* \mid \theta_{-k}, x)} \\
  &= \frac{\pi(\theta_k^* \mid \theta_{-k}, x)\,\pi(\theta_{-k} \mid x)
            \cdot \pi(\theta_k^{(t)} \mid \theta_{-k}, x)}
          {\pi(\theta_k^{(t)} \mid \theta_{-k}, x)\,\pi(\theta_{-k} \mid x)
            \cdot \pi(\theta_k^* \mid \theta_{-k}, x)}
  = 1.
\end{align*}
\end{proof}

\begin{theoreme}[Convergence de l'\'echantillonneur de Gibbs]
Sous des conditions de r\'egularit\'e (positivit\'e du posterior sur
tout l'espace), l'\'echantillonneur de Gibbs converge vers la loi
cible $\pi(\theta \mid x)$.
\end{theoreme}

% ------------------------------------------------------------
\section{Exemples fondamentaux}
% ------------------------------------------------------------

\subsection{Mod\`ele normal bivari\'e}

\begin{exemple}
Soit $(\theta_1, \theta_2) \sim \mathcal{N}_2(\boldsymbol{\mu}, \boldsymbol{\Sigma})$
avec $\boldsymbol{\Sigma} = \begin{pmatrix} 1 & \rho \\ \rho & 1 \end{pmatrix}$.
Les conditionnelles compl\`etes sont~:
\begin{align*}
  \theta_1 \mid \theta_2 &\sim \mathcal{N}(\mu_1 + \rho(\theta_2 - \mu_2),\; 1 - \rho^2), \\
  \theta_2 \mid \theta_1 &\sim \mathcal{N}(\mu_2 + \rho(\theta_1 - \mu_1),\; 1 - \rho^2).
\end{align*}
\end{exemple}

\begin{attention}[title=\textbf{Forte corr\'elation}]
Quand $\rho$ est proche de $\pm 1$, le Gibbs sampler m\'elange lentement
car les mises \`a jour par composante ne peuvent pas explorer
efficacement la distribution allong\'ee. On pr\'ef\`ere alors HMC/NUTS
ou une reparam\'etrisation.
\end{attention}

\subsection{Mod\`ele normal avec moyenne et variance inconnues}

\begin{exemple}
Soit $X_1, \ldots, X_n \overset{\text{iid}}{\sim} \mathcal{N}(\mu, \sigma^2)$
avec prior $\mu \sim \mathcal{N}(\mu_0, \tau_0^2)$ et
$\sigma^2 \sim \text{IG}(\alpha_0, \beta_0)$. Les conditionnelles sont~:
\begin{align*}
  \mu \mid \sigma^2, x &\sim \mathcal{N}\!\left(
  \frac{\mu_0/\tau_0^2 + n\bar{x}/\sigma^2}{1/\tau_0^2 + n/\sigma^2},\;
  \frac{1}{1/\tau_0^2 + n/\sigma^2}\right), \\
  \sigma^2 \mid \mu, x &\sim \text{IG}\!\left(
  \alpha_0 + \frac{n}{2},\;
  \beta_0 + \frac{1}{2}\sum_{i=1}^n (x_i - \mu)^2\right).
\end{align*}
\end{exemple}

\subsection{Mod\`ele de m\'elange}

\begin{exemple}
\textbf{M\'elange de $K$ gaussiennes.} Soient les donn\'ees
$x_1, \ldots, x_n$ avec variables latentes $z_i \in \{1, \ldots, K\}$
et param\`etres $(\boldsymbol{\pi}, \boldsymbol{\mu}, \boldsymbol{\sigma}^2)$.
Les conditionnelles sont~:

\textbf{Allocation}~:
$\PP(z_i = k \mid \cdot) \propto \pi_k\,\mathcal{N}(x_i \mid \mu_k, \sigma_k^2)$.

\textbf{Proportions}~:
$\boldsymbol{\pi} \mid \cdot \sim \text{Dir}(\alpha_1 + n_1, \ldots, \alpha_K + n_K)$,
o\`u $n_k = \#\{i : z_i = k\}$.

\textbf{Moyennes}~:
$\mu_k \mid \cdot \sim \mathcal{N}\!\left(\frac{\mu_0/\tau^2 + \sum_{z_i=k} x_i/\sigma_k^2}{1/\tau^2 + n_k/\sigma_k^2},\;
\frac{1}{1/\tau^2 + n_k/\sigma_k^2}\right)$.

\textbf{Variances}~:
$\sigma_k^2 \mid \cdot \sim \text{IG}\!\left(a + n_k/2,\;
b + \frac{1}{2}\sum_{z_i=k}(x_i - \mu_k)^2\right)$.
\end{exemple}

% ------------------------------------------------------------
\section{Variantes de Gibbs}
% ------------------------------------------------------------

\begin{definition}[Gibbs par blocs]
Au lieu de mettre \`a jour chaque composante s\'epar\'ement, le
\textbf{Gibbs par blocs} met \`a jour des groupes de composantes
simultan\'ement. Cela am\'eliore le m\'elange quand les composantes
d'un bloc sont fortement corr\'el\'ees.
\end{definition}

\begin{definition}[Collapsed Gibbs / Rao--Blackwellisation]
Le \textbf{collapsed Gibbs} marginalise analytiquement certains
param\`etres avant l'\'echantillonnage, r\'eduisant la dimension et
am\'eliorant l'efficacit\'e.
\end{definition}

\begin{theoreme}[Rao--Blackwellisation]
Si $g(\theta)$ est une fonction d'int\'er\^et et qu'on peut calculer
$\E[g(\theta_1) \mid \theta_2, x]$ analytiquement, alors l'estimateur
Rao--Blackwellis\'e~:
\[
  \hat{g}_{\text{RB}} = \frac{1}{T}\sum_{t=1}^T
  \E[g(\theta_1) \mid \theta_2^{(t)}, x]
\]
a une variance inf\'erieure ou \'egale \`a $\hat{g} = \frac{1}{T}\sum_t g(\theta_1^{(t)})$.
\end{theoreme}

% ------------------------------------------------------------
\section{Augmentation de donn\'ees}
% ------------------------------------------------------------

\begin{definition}[Augmentation de donn\'ees]
L'\textbf{augmentation de donn\'ees} introduit des variables latentes
$z$ pour faciliter l'\'echantillonnage. On alterne~:
\begin{enumerate}
  \item Tirer $z \mid \theta, x$ (\'etape I).
  \item Tirer $\theta \mid z, x$ (\'etape P).
\end{enumerate}
\end{definition}

\begin{exemple}
\textbf{Mod\`ele probit.} Pour $Y_i \in \{0, 1\}$ avec
$\PP(Y_i = 1) = \Phi(X_i^\top \beta)$, on introduit
$Z_i \sim \mathcal{N}(X_i^\top \beta, 1)$ et $Y_i = \mathbf{1}(Z_i > 0)$.
Conditionnellement aux $Z_i$, $\beta$ a un posterior gaussien. L'\'echantillonnage
de $Z_i$ se fait par troncature.
\end{exemple}

% ------------------------------------------------------------
\section{Impl\'ementation Python}
% ------------------------------------------------------------

\begin{codeblock}[title=\textbf{Gibbs pour le mod\`ele normal}]
\begin{minted}{python}
import numpy as np
import matplotlib.pyplot as plt
from scipy import stats

# Données
np.random.seed(42)
n = 50
mu_true, sigma_true = 5.0, 2.0
data = np.random.normal(mu_true, sigma_true, n)
xbar = data.mean()
S2 = ((data - data.mean())**2).sum()

# Hyperparamètres
mu_0, tau_0 = 0, 10
alpha_0, beta_0 = 1, 1

# Gibbs sampler
T = 10000
mu_samples = np.zeros(T)
sigma2_samples = np.zeros(T)
mu_samples[0] = 0
sigma2_samples[0] = 1

for t in range(1, T):
    # Tirer mu | sigma^2, x
    prec_post = 1/tau_0**2 + n/sigma2_samples[t-1]
    mu_post = (mu_0/tau_0**2 + n*xbar/sigma2_samples[t-1]) \
              / prec_post
    mu_samples[t] = np.random.normal(mu_post,
                                      1/np.sqrt(prec_post))

    # Tirer sigma^2 | mu, x
    alpha_post = alpha_0 + n/2
    beta_post = beta_0 + 0.5 * np.sum(
        (data - mu_samples[t])**2)
    sigma2_samples[t] = 1 / np.random.gamma(alpha_post,
                                              1/beta_post)

burn_in = 1000
print(f"mu: {mu_samples[burn_in:].mean():.3f} "
      f"(vrai: {mu_true})")
print(f"sigma: {np.sqrt(sigma2_samples[burn_in:]).mean():.3f} "
      f"(vrai: {sigma_true})")
\end{minted}
\end{codeblock}

\begin{codeblock}[title=\textbf{Gibbs pour m\'elange de gaussiennes}]
\begin{minted}{python}
# Mélange de 2 gaussiennes
np.random.seed(42)
n = 200
z_true = np.random.choice([0, 1], size=n, p=[0.4, 0.6])
data_mix = np.where(z_true == 0,
                     np.random.normal(0, 1, n),
                     np.random.normal(5, 1.5, n))

K = 2
T = 5000

# Stockage
pi_samples = np.zeros((T, K))
mu_samples_mix = np.zeros((T, K))
sigma2_samples_mix = np.zeros((T, K))
z_samples = np.zeros((T, n), dtype=int)

# Initialisation
mu_samples_mix[0] = [-1, 6]
sigma2_samples_mix[0] = [1, 1]
pi_samples[0] = [0.5, 0.5]

for t in range(1, T):
    # Étape E: tirer z_i
    for i in range(n):
        probs = np.array([
            pi_samples[t-1, k] *
            stats.norm.pdf(data_mix[i],
                           mu_samples_mix[t-1, k],
                           np.sqrt(sigma2_samples_mix[t-1, k]))
            for k in range(K)])
        probs /= probs.sum()
        z_samples[t, i] = np.random.choice(K, p=probs)

    # Étape M: tirer paramètres
    for k in range(K):
        idx = z_samples[t] == k
        nk = idx.sum()
        # Proportions
        pi_samples[t, k] = nk  # sera normalisé
        if nk > 0:
            xk = data_mix[idx]
            # mu_k
            prec = 1/100 + nk/sigma2_samples_mix[t-1, k]
            m = (nk * xk.mean() / sigma2_samples_mix[t-1, k]) / prec
            mu_samples_mix[t, k] = np.random.normal(m, 1/np.sqrt(prec))
            # sigma^2_k
            a = 1 + nk/2
            b = 1 + 0.5 * np.sum((xk - mu_samples_mix[t, k])**2)
            sigma2_samples_mix[t, k] = 1/np.random.gamma(a, 1/b)
        else:
            mu_samples_mix[t, k] = mu_samples_mix[t-1, k]
            sigma2_samples_mix[t, k] = sigma2_samples_mix[t-1, k]

    # Dirichlet pour les proportions
    counts = np.array([(z_samples[t]==k).sum() for k in range(K)])
    pi_samples[t] = np.random.dirichlet(1 + counts)

burn = 1000
print("mu estimates:", mu_samples_mix[burn:].mean(axis=0))
print("pi estimates:", pi_samples[burn:].mean(axis=0))
\end{minted}
\end{codeblock}

\begin{codeblock}[title=\textbf{Diagnostics et visualisation}]
\begin{minted}{python}
fig, axes = plt.subplots(2, 2, figsize=(12, 8))

# Trace plots
axes[0, 0].plot(mu_samples[:3000])
axes[0, 0].set_title(r"Trace de $\mu$")
axes[0, 1].plot(sigma2_samples[:3000])
axes[0, 1].set_title(r"Trace de $\sigma^2$")

# Joint posterior
axes[1, 0].scatter(mu_samples[burn_in::10],
                    sigma2_samples[burn_in::10],
                    alpha=0.1, s=1)
axes[1, 0].set_xlabel(r"$\mu$")
axes[1, 0].set_ylabel(r"$\sigma^2$")
axes[1, 0].set_title("Posterior joint")

# Marginal de sigma
axes[1, 1].hist(np.sqrt(sigma2_samples[burn_in:]),
                bins=50, density=True, alpha=0.5)
axes[1, 1].axvline(sigma_true, color="red",
                    linestyle="--")
axes[1, 1].set_title(r"Marginal de $\sigma$")

plt.tight_layout()
plt.savefig("gibbs_diagnostics.pdf")
\end{minted}
\end{codeblock}

% ------------------------------------------------------------
\section{Exercices}
% ------------------------------------------------------------

\begin{exercice}
D\'eriver les conditionnelles compl\`etes pour le mod\`ele de r\'egression
lin\'eaire bay\'esien $y = X\beta + \varepsilon$ avec
$\beta \sim \mathcal{N}(0, \sigma_\beta^2 I)$ et
$\sigma^2 \sim \text{IG}(\alpha, \beta)$.
\end{exercice}

\begin{exercice}
Impl\'ementer l'augmentation de donn\'ees d'Albert et Chib (1993) pour
le mod\`ele probit et comparer avec PyMC.
\end{exercice}

\begin{exercice}
Pour le Gibbs bivari\'e gaussien, montrer analytiquement que
l'autocorr\'elation d'ordre $k$ du Gibbs sampler est $\rho^{2k}$ et
calculer l'ESS en fonction de $\rho$.
\end{exercice}

\begin{exercice}
Comparer Gibbs standard et Gibbs par blocs sur le mod\`ele normal
multivari\'e avec diff\'erentes structures de corr\'elation.
\end{exercice}

\begin{exercice}
\textbf{(Avanc\'e)} Impl\'ementer le collapsed Gibbs pour le mod\`ele
de m\'elange en marginalisant les proportions $\pi_k$. Comparer l'ESS
avec le Gibbs standard.
\end{exercice}
